According to documents from **July 7, 2026**, the intended model is the projected analytical map

[
x
\longmapsto
(U_x,D_x,P_x,\mathcal A,\Psi)
\longmapsto
(c,K,Q,d,M,D^{\rm self},A,h)
\longmapsto
\sigma,
]

with no production Green–Kubo time integral and no Einstein–Helfand slope in the prediction path. The same documents also state that a recipe containing only names and concentrations is not sufficient unless it includes the physical parameter map needed to construct (U_x,D_x,P_x,A_i,\psi_\alpha). 

Below is the full implementation-grade mathematical model.

---

# Complete projected analytical electrolyte conductivity model

## 0. Output

Given one electrolyte composition (x), return one scalar conductivity

[
\sigma(x)\quad [\mathrm{S,m^{-1}}],
]

or equivalently

[
\sigma_{\mathrm{mS/cm}}(x)=10,\sigma(x).
]

The production model is:

[
\boxed{
x
\rightarrow
(U_x,D_x,P_x,\mathcal A,\Psi)
\rightarrow
(c,K,Q,d,M,D^{\rm self},A,h)
\rightarrow
\sigma
}
]

where:

[
U_x(q)
]

is the molar free energy,

[
D_x(q)
]

is the microscopic or reduced mobility/diffusion tensor,

[
P_x(q)
]

is the unwrapped charge-polarization observable,

[
\mathcal A={A_i}_{i=1}^N
]

is the finite partition into transport basins,

[
\Psi=(\psi_1,\ldots,\psi_m)
]

is the finite memory-coordinate basis, and

[
(c,K,Q,d,M,D^{\rm self},A,h)
]

are the projected conductivity primitives.

No production term is allowed outside these primitives.

---

# 1. Composition input

The full composition is not just a recipe. It is:

[
x=(r,\mathcal L),
]

where

[
r=
\left(
T,,
{C_a}*{a\in\mathcal C},,
{\phi_s}*{s\in\mathcal S}
\right)
]

and (\mathcal L) is the physical library.

Definitions:

[
T>0
]

is temperature in K.

[
\mathcal C
]

is the set of conserved components.

Examples:

[
\mathcal C=
{\mathrm{Li},\mathrm{PF6},\mathrm{FSI},\mathrm{EC},\mathrm{DMC},\mathrm{FEC},\ldots}.
]

[
C_a\quad [\mathrm{mol,m^{-3}}]
]

is the analytical total concentration of conserved component (a).

[
\phi_s
]

is the liquid fraction of solvent/additive species (s).

The physical library is

[
\mathcal L=
\left(
{\mathcal T_s}*{s\in\mathcal S*{\rm mol}},
{\mathcal Q_s}*{s\in\mathcal S*{\rm mol}},
{\Theta_{st}}_{s,t},
\mathcal B,
\mathcal M,
\mathcal R,
\mathcal \Psi
\right).
]

Here:

[
\mathcal T_s
]

is topology of species (s),

[
\mathcal Q_s
]

is its charge-site model,

[
\Theta_{st}
]

are interaction parameters between species (s,t),

[
\mathcal B
]

is the basin-definition library,

[
\mathcal M
]

is the mobility/friction library,

[
\mathcal R
]

is the transition-surface / reaction-coordinate library,

[
\mathcal\Psi
]

is the Mori memory-coordinate library.

A recipe-only object such as

[
{\mathrm{EC}:,0.3,\ \mathrm{DMC}:,0.7,\ \mathrm{LiPF6}:,1.0\mathrm{M}}
]

does not determine the model unless it also determines (\mathcal L). Otherwise many different (U_x,D_x,P_x,A_i,\Psi) are compatible with the same recipe.

---

# 2. Non-identifiability of recipe-only input

Suppose two physical libraries (\mathcal L_1,\mathcal L_2) have the same species names and same concentrations but different microscopic mobility tensors:

[
D_x^{(2)}(q)=\alpha D_x^{(1)}(q),
\qquad
\alpha\neq 1.
]

Keeping (U_x,P_x,A_i,\Psi) fixed, the self-current tensors scale as

[
D_{i,2}^{\rm self}=\alpha D_{i,1}^{\rm self},
]

and, for overdamped Smoluchowski dynamics, the capacity fluxes scale as

[
K_{ij}^{(2)}=\alpha K_{ij}^{(1)}
]

because

[
K_{ij}=
\int
\nabla q_{ij}(q)^\top D_x(q)\nabla q_{ij}(q),
\rho_x(\mathrm dq).
]

Therefore the predicted conductivity generally changes:

[
\sigma(r,\mathcal L_2)\neq \sigma(r,\mathcal L_1).
]

So the complete input is

[
x=(r,\mathcal L),
]

not (r) alone.

This is the reason a descriptor-only recipe path is a surrogate unless it really constructs the physical generator. The retrieved notes state the same boundary: descriptor vectors must define (U_x), boundaries (A_i), reactive fluxes (K_{ij}), and moments (d_{ij}), not arbitrary fitted primitives. 

---

# 3. Configuration space

Let

[
q\in \Omega_x\subseteq \mathbb R^d
]

be the coordinate vector.

For atomistic coordinates,

[
q=
(r_{1,1},\ldots,r_{1,n_1},r_{2,1},\ldots,r_{N_{\rm mol},n_{N_{\rm mol}}}),
]

with

[
r_{m,u}\in\mathbb R^3.
]

Here (m) indexes molecules and (u) indexes sites in molecule (m).

The full coordinate dimension is

[
d=3\sum_{m=1}^{N_{\rm mol}} n_{s(m)}.
]

For a reduced coordinate model,

[
q=
\left(
r_{\rm Li-A},
s_{\rm pair},
n_{\rm Li-solv},
n_{\rm Li-lig},
\omega_{\rm A},
n_{\rm cluster},
\rho_{\rm local},
p_{\rm atm},
\ldots
\right).
]

Every coordinate must have:

1. a unit,
2. a physical definition,
3. a gradient,
4. a mobility entry in (D_x(q)),
5. a role in (U_x(q)), (P_x(q)), a basin (A_i), or a memory coordinate (\psi_\alpha(q)).

---

# 4. Species topology and charge sites

For each species (s), define site set

[
\mathcal I_s={1,\ldots,n_s}.
]

Each site (u\in \mathcal I_s) has:

[
m_{su}>0
]

mass,

[
a_{su}>0
]

hydrodynamic radius,

[
\sigma_{su}>0,\qquad \epsilon_{su}\ge0
]

short-range nonbonded parameters,

[
z_{su}\in\mathbb R
]

formal or partial charge number,

[
\alpha_{su}\ge0
]

polarizability, if used.

The topology contains bonded edges

[
\mathcal E_s^{\rm bond}\subset\mathcal I_s\times\mathcal I_s,
]

angles

[
\mathcal E_s^{\rm angle}\subset\mathcal I_s^3,
]

torsions

[
\mathcal E_s^{\rm torsion}\subset\mathcal I_s^4.
]

The charge of site ((m,u)) is

[
z_{m,u}=z_{s(m),u}.
]

The SI charge is

[
q_{m,u}^{\rm SI}=e z_{m,u}.
]

---

# 5. Molecule counts from composition

For an explicit simulation cell of volume (V_{\rm cell}), concentration (C_a) gives particle count

[
N_a=\operatorname{round}(C_a N_A V_{\rm cell}).
]

Here

[
N_A
]

is Avogadro’s number.

For solvent/additive species given by volume fraction (\phi_s), molar mass (M_s), and density (\rho_s), pure number density is

[
n_s^{\rm pure}=\frac{\rho_s}{M_s}N_A.
]

The target count is

[
N_s=
\operatorname{round}\left(\phi_s n_s^{\rm pure}V_{\rm cell}\right).
]

If volume is not fixed and both concentrations and liquid fractions are specified, solve the volume-filling equation

[
V_{\rm cell}
============

\sum_s \frac{N_sM_s}{\rho_sN_A}
+
\sum_a \bar v_a N_a,
]

where (\bar v_a) is the partial molar volume of component (a).

Fixed-point algorithm:

1. Initialize (V_{\rm cell}^{(0)}).
2. Compute (N_a^{(k)}=\operatorname{round}(C_aN_AV_{\rm cell}^{(k)})).
3. Compute (N_s^{(k)}=\operatorname{round}(\phi_s n_s^{\rm pure}V_{\rm cell}^{(k)})).
4. Update

[
V_{\rm cell}^{(k+1)}
====================

\sum_s \frac{N_s^{(k)}M_s}{\rho_sN_A}
+
\sum_a\bar v_aN_a^{(k)}.
]

5. Stop when

[
\frac{|V_{\rm cell}^{(k+1)}-V_{\rm cell}^{(k)}|}
{V_{\rm cell}^{(k)}}<\tau_V.
]

---

# 6. Charge polarization

The formal-charge polarization is

[
P_x(q)
======

\sum_m\sum_{u\in\mathcal I_{s(m)}}
z_{m,u}R_{m,u}^{\rm unwrap}(q)
\in\mathbb R^3.
]

Here (R_{m,u}^{\rm unwrap}) is the unwrapped site position.

Component form:

[
P_{x,a}(q)
==========

\sum_m\sum_u z_{m,u}R_{m,u,a}^{\rm unwrap}(q),
\qquad
a\in{x,y,z}.
]

The gradient is

[
G_P(q)=\nabla_qP_x(q)\in\mathbb R^{3\times d}.
]

For coordinate (r_{m,u,\ell}),

[
\frac{\partial P_{x,a}}{\partial r_{m,u,\ell}}
==============================================

z_{m,u}\delta_{a\ell}.
]

If reduced coordinates (\xi=\Xi(q)) are used, then

[
\nabla_\xi P_x
==============

\nabla_qP_x
\left(\frac{\partial \Xi}{\partial q}\right)^#.
]

Here (B^#) means Moore–Penrose pseudoinverse.

---

# 7. Pseudoinverse operation (B^#)

For symmetric positive semidefinite (B):

1. Eigendecompose

[
B=V\Lambda V^\top,
]

where

[
\Lambda=\operatorname{diag}(\lambda_1,\ldots,\lambda_m),
\qquad
\lambda_r\ge0.
]

2. Choose tolerance

[
\tau_#=10^{-12}\max_r\lambda_r.
]

3. Define

[
\lambda_r^#=
\begin{cases}
1/\lambda_r, & \lambda_r>\tau_#,\
0, & \lambda_r\le\tau_#.
\end{cases}
]

4. Return

[
B^#=V\operatorname{diag}(\lambda_1^#,\ldots,\lambda_m^#)V^\top.
]

For a current coupling vector (y), require the nullspace consistency condition

[
|(I-BB^#)y|
\le
\tau_{\rm null}|y|,
]

for example

[
\tau_{\rm null}=10^{-8}.
]

If this fails, the memory basis is ill-posed because current couples to a zero-energy mode.

---

# 8. Free energy (U_x(q))

The molar free energy is

[
U_x:\Omega_x\to\mathbb R
\qquad [\mathrm{J,mol^{-1}}].
]

Use

[
U_x(q)
======

U_{\rm intra}(q)
+
U_{\rm LJ}(q)
+
U_{\rm Coul}(q)
+
U_{\rm pol}^{\rm eff}(q)
+
U_{\rm coord}(q)
+
U_{\rm solv}(q)
+
U_{\rm pack}(q)
+
U_{\rm atm}(q).
]

## 8.1 Bonded terms

Bond energy:

[
U_{\rm bond}(q)
===============

N_A
\sum_m
\sum_{(u,v)\in\mathcal E_{s(m)}^{\rm bond}}
\frac12 k_{uv}^{\rm bond}
\left(
|r_{m,u}-r_{m,v}|-r_{uv}^0
\right)^2.
]

Angle energy:

[
U_{\rm angle}(q)
================

N_A
\sum_m
\sum_{(u,v,w)\in\mathcal E_{s(m)}^{\rm angle}}
\frac12 k_{uvw}^{\rm angle}
\left(
\theta_{uvw}(q)-\theta_{uvw}^0
\right)^2.
]

Torsion energy:

[
U_{\rm torsion}(q)
==================

N_A
\sum_m
\sum_{(u,v,w,z)\in\mathcal E_{s(m)}^{\rm torsion}}
\sum_{n=1}^{n_{\max}}
\frac12 V_n
\left[
1+\cos(n\varphi_{uv wz}(q)-\gamma_n)
\right].
]

Then

[
U_{\rm intra}=U_{\rm bond}+U_{\rm angle}+U_{\rm torsion}.
]

## 8.2 Lennard–Jones / repulsion–dispersion

For site pair (a,b),

[
r_{ab}=|r_a-r_b|_{\rm PBC}.
]

Mixing rules:

[
\sigma_{ab}=\frac{\sigma_a+\sigma_b}{2},
\qquad
\epsilon_{ab}=\sqrt{\epsilon_a\epsilon_b}.
]

Potential:

[
U_{\rm LJ}(q)
=============

N_A
\sum_{a<b}
4\epsilon_{ab}
\left[
\left(\frac{\sigma_{ab}}{r_{ab}}\right)^{12}
--------------------------------------------

\left(\frac{\sigma_{ab}}{r_{ab}}\right)^6
\right]
f_{\rm cut}(r_{ab}).
]

Switching function:

[
f_{\rm cut}(r)=
\begin{cases}
1, & r\le r_{\rm on},\
1-10s^3+15s^4-6s^5, & r_{\rm on}<r<r_{\rm off},\
0, & r\ge r_{\rm off},
\end{cases}
]

with

[
s=\frac{r-r_{\rm on}}{r_{\rm off}-r_{\rm on}}.
]

## 8.3 Electrostatics

Use Ewald electrostatics:

[
U_{\rm Coul}(q)
===============

N_A
\frac{1}{4\pi\epsilon_0\epsilon_\infty}
\left[
U_{\rm real}
+
U_{\rm reciprocal}
+
U_{\rm self}
+
U_{\rm surface}
\right].
]

Real-space term:

[
U_{\rm real}
============

\frac12
\sum_{a\ne b}
q_a^{\rm SI}q_b^{\rm SI}
\frac{\operatorname{erfc}(\alpha r_{ab})}{r_{ab}}.
]

Reciprocal-space term:

[
U_{\rm reciprocal}
==================

\frac{2\pi}{V_{\rm cell}}
\sum_{k\ne0}
\frac{\exp[-k^2/(4\alpha^2)]}{k^2}
\left|
\sum_a q_a^{\rm SI}\exp(ik\cdot r_a)
\right|^2.
]

Self term:

[
U_{\rm self}
============

-\frac{\alpha}{\sqrt\pi}
\sum_a(q_a^{\rm SI})^2.
]

For conducting boundary condition:

[
U_{\rm surface}=0.
]

## 8.4 Polarization

If polarization is used, define induced dipoles (\mu_a), polarizabilities (\alpha_a), electric fields (E_a(q)), and dipole interaction tensors (T_{ab}):

[
U_{\rm pol}(q,{\mu_a})
======================

N_A
\left[
\sum_a\frac12\mu_a^\top\alpha_a^{-1}\mu_a
-----------------------------------------

\sum_a\mu_a^\top E_a(q)
+
\frac12\sum_{a\ne b}\mu_a^\top T_{ab}\mu_b
\right].
]

Minimize:

[
{\mu_a^\star(q)}
================

\arg\min_{{\mu_a}}
U_{\rm pol}(q,{\mu_a}).
]

Then

[
U_{\rm pol}^{\rm eff}(q)=U_{\rm pol}(q,{\mu_a^\star(q)}).
]

If polarization is not used:

[
U_{\rm pol}^{\rm eff}(q)=0.
]

## 8.5 Coordination and solvation

Smooth coordination switch:

[
s_{aL}(r)=\frac{1}{1+(r/r_{aL}^0)^{p_{aL}}}.
]

Coordination number:

[
n_{aL}(q)=\sum_{b\in L}s_{aL}(r_{ab}).
]

Coordination energy:

[
U_{\rm coord}(q)
================

N_A
\sum_{a,L}
\Delta g_{aL}^{\rm coord}
,n_{aL}(q).
]

Solvation-shell competition:

[
U_{\rm solv}(q)
===============

N_A
\sum_a
F_a^{\rm solv}(n_{aL_1},n_{aL_2},\ldots).
]

A concrete quadratic form is

[
F_a^{\rm solv}(n_1,\ldots,n_m)
==============================

\frac12
\sum_{\ell,\ell'}
K_{\ell\ell'}^a
(n_\ell-\bar n_\ell^a)
(n_{\ell'}-\bar n_{\ell'}^a).
]

## 8.6 Packing / free volume

Let local occupied volume fraction be

[
\varphi(q)=\sum_a v_a W_a(q),
]

where (v_a) is molecular volume and (W_a(q)) is a local kernel.

Use Carnahan–Starling-like packing penalty:

[
U_{\rm pack}(q)
===============

N_A RT
\sum_\ell
w_\ell^{\rm grid}
\frac{
\varphi_\ell(4-3\varphi_\ell)
}{
(1-\varphi_\ell)^2
}.
]

---

# 9. Equilibrium measure

Define

[
\beta_{\rm mol}=\frac{1}{RT}.
]

The equilibrium measure is

[
\mu_x(\mathrm dq)=Z_x^{-1}\exp[-\beta_{\rm mol}U_x(q)],\mathrm dq.
]

---

# 10. Mobility / diffusion tensor (D_x(q))

Define resistance matrix

[
R_x(q)\in\mathbb R^{d\times d}.
]

Decompose

[
R_x(q)
======

R_{\rm Stokes}(q)
+
R_{\rm hydro}(q)
+
R_{\rm constraint}(q)
+
R_{\rm atmosphere}(q)
+
R_{\rm cage}(q)
+
R_{\rm microvisc}(q).
]

Then

[
D_x(q)=k_BT,R_x(q)^#.
]

Require

[
D_x(q)=D_x(q)^\top,
\qquad
D_x(q)\succeq0.
]

## 10.1 Local viscosity

[
\eta(q)
=======

\eta_{\rm bulk}(x)
f_{\rm salt}(I(q))
f_{\rm additive}(q)
f_{\rm free\ volume}(\varphi(q)).
]

Use

[
f_{\rm salt}(I)=1+A_\eta\sqrt I+B_\eta I,
]

[
f_{\rm additive}(q)=1+\sum_L a_L n_L(q),
]

[
f_{\rm free\ volume}(\varphi)
=============================

\left(1-\frac{\varphi}{\varphi_{\max}}\right)^{-\gamma}.
]

## 10.2 Stokes resistance

For site (a):

[
\zeta_a(q)=6\pi\eta(q)a_a.
]

Then

[
R_{{\rm Stokes},ab}(q)=\zeta_a(q)I_3\delta_{ab}.
]

## 10.3 Hydrodynamic cross mobility

For (a\ne b), use Rotne–Prager–Yamakawa mobility when desired:

[
H_{ab}(q)
=========

\frac{1}{8\pi\eta r_{ab}}
\left[
I_3+\hat r_{ab}\hat r_{ab}^\top
+
\frac{2(a_a^2+a_b^2)}{3r_{ab}^2}
\left(
I_3-3\hat r_{ab}\hat r_{ab}^\top
\right)
\right].
]

For (a=b):

[
H_{aa}=\frac{1}{6\pi\eta a_a}I_3.
]

Then

[
D_{\rm hydro}=k_BT,H.
]

## 10.4 Constraint resistance

For constraint mode (m), with direction (s_m(q)), lifetime (\tau_m(q)), and length (\ell_m(q)),

[
R_{\rm constraint}^{(m)}(q)
===========================

k_BT\frac{\tau_m(q)}{\ell_m(q)^2}
s_m(q)s_m(q)^\top.
]

Then

[
R_{\rm constraint}(q)=\sum_mR_{\rm constraint}^{(m)}(q).
]

## 10.5 Ion-atmosphere resistance

Define charge density mode

[
\rho_z(k;q)=\sum_a z_a\exp(ik\cdot r_a).
]

Define inverse Debye length

[
\kappa^2(q)
===========

\frac{F^2}{\epsilon_0\epsilon_rRT}
\sum_i z_i^2 c_i(q).
]

Let

[
g_k(q)=\nabla_q\rho_z(k;q).
]

Then an admissible atmosphere resistance is

[
R_{\rm atmosphere}(q)
=====================

\sum_{k\in\mathcal K}
\Gamma(k;\kappa,\eta,\epsilon_r)
g_k(q)g_k(q)^\top.
]

A simple kernel is

[
\Gamma(k;\kappa,\eta,\epsilon_r)
================================

\Gamma_0
\frac{k^2}{(k^2+\kappa^2)^2}
\eta.
]

This contributes to (D_x(q)), (A), and (h). It is not a final multiplicative conductivity correction.

---

# 11. Reversible generator (L_x)

The overdamped Smoluchowski generator is

[
L_x f(q)
========

\frac{1}{\mu_x(q)}
\nabla_q\cdot
\left[
\mu_x(q)D_x(q)\nabla_qf(q)
\right].
]

Equivalently,

[
L_x f
=====

D_x:\nabla_q\nabla_q f
+
\left(
\nabla_q\cdot D_x
-----------------

\beta_{\rm mol}D_x\nabla_qU_x
\right)\cdot\nabla_qf.
]

---

# 12. Reduced generator, if used

If reduced coordinates are used,

[
\xi=\Xi(q)\in\mathbb R^m,
]

then the reduced PMF must be derived from the full generator:

[
U_{\rm red}(\xi)
================

-RT\log
\int_{\Omega_x}
\delta(\Xi(q)-\xi)
\exp[-U_x(q)/(RT)],\mathrm dq.
]

The reduced mobility is

[
D_{\rm red}(\xi)
================

\mathbb E_{\mu_x}
\left[
\nabla_q\Xi(q)D_x(q)\nabla_q\Xi(q)^\top
\mid
\Xi(q)=\xi
\right].
]

The reduced charge polarization is

[
P_{\rm red}(\xi)
================

\mathbb E_{\mu_x}
\left[
P_x(q)\mid \Xi(q)=\xi
\right].
]

A reduced model is first-principles only if (U_{\rm red},D_{\rm red},P_{\rm red}), basins, and memory coordinates are conditional projections from the full generator.

---

# 13. Basin definitions (A_i)

A basin is a measurable subset of (\Omega_x):

[
A_i\subset\Omega_x.
]

Basins are disjoint:

[
A_i\cap A_j=\varnothing,\quad i\ne j.
]

Every coordinate is assigned to one basin:

[
s(q)=i\quad\Longleftrightarrow\quad q\in A_i.
]

Each basin has a stoichiometry vector

[
\nu_i=(\nu_{i,a})_{a\in\mathcal C}.
]

Here

[
\nu_{i,a}
]

is the number of conserved component (a) in motif (i).

## 13.1 Pair basins

For Li (\ell) and anion (a):

[
r_{\ell a}=|r_\ell-r_a|.
]

Contact ion pair:

[
A_{\rm CIP}={q:r_{\ell a}<r_{\rm CIP}}.
]

Solvent-separated ion pair:

[
A_{\rm SSIP}={q:r_{\rm CIP}\le r_{\ell a}<r_{\rm SSIP}}.
]

Free basin:

[
A_{\rm free}=
{q:r_{\ell a}\ge r_{\rm SSIP}\ \text{for all nearby anions}}.
]

## 13.2 Ligand basins

Li-ligand coordination:

[
n_{\ell L}(q)=\sum_{b\in L}s_{\ell b}(r_{\ell b}).
]

Li-ligand basin:

[
A_{\rm LiL}={q:n_{\ell L}\ge n_{\rm lig}^{\rm cut}}.
]

Ligand-separated SSIP:

[
A_{\rm LiL\cdots A}
===================

A_{\rm SSIP}\cap A_{\rm LiL}.
]

## 13.3 Cluster basins

Build graph

[
G(q)=(V,E(q)).
]

Vertices are charged species. Edge ((a,b)\in E(q)) exists when

[
r_{ab}<r_{\rm pair}^{\rm cut}.
]

Clusters are connected components of (G(q)).

Cluster basin labels include:

[
A_{\rm free},\quad A_{\rm CIP},\quad A_{\rm SSIP},\quad A_{\rm aggregate},\quad A_{\rm bridge}.
]

## 13.4 Orientation basins

For anion orientation vector (u_a(q)),

[
\omega_a(q)=u_a(q)\cdot \hat r_{\ell a},
]

where

[
\hat r_{\ell a}=\frac{r_a-r_\ell}{|r_a-r_\ell|}.
]

Bin:

[
A_{\omega,k}={q:\omega_k^-\le\omega_a(q)<\omega_k^+}.
]

## 13.5 Joint basis

The final basin is a joint bin:

[
A_i=
A_{\rm pair}
\cap
A_{\rm ligand}
\cap
A_{\rm cluster}
\cap
A_{\rm orientation}
\cap
A_{\rm cage}
\cap
A_{\rm atmosphere}.
]

The retained set of nonempty joint bins is

[
\mathcal A={A_1,\ldots,A_N}.
]

---

# 14. Memory coordinates (\Psi)

Memory coordinates are functions

[
\psi_\alpha:\Omega_x\to\mathbb R.
]

Collect:

[
\Psi(q)=(\psi_1(q),\ldots,\psi_m(q))^\top.
]

Examples:

## 14.1 Charge-density modes

[
\psi_{k,c}(q)=\sum_a z_a\cos(k\cdot r_a),
]

[
\psi_{k,s}(q)=\sum_a z_a\sin(k\cdot r_a).
]

## 14.2 Cage coordinate

Let (R_{\rm cage}) be the solvent cage center:

[
R_{\rm cage}(q)=
\frac{\sum_{b\in\mathcal N_\ell(q)} w_b r_b}
{\sum_{b\in\mathcal N_\ell(q)} w_b}.
]

Then

[
\psi_{\rm cage}(q)=(r_\ell-R_{\rm cage})\cdot e_{\rm cage}.
]

## 14.3 Partner residence

For current partner (a^\star(q)):

[
\psi_{\rm partner}(q)=s_{\ell a^\star}(q).
]

## 14.4 Ligand residence

[
\psi_{\rm ligand}(q)=n_{\ell L}(q)-\bar n_{\ell L}.
]

## 14.5 Anion orientation

[
\psi_{\rm orient}(q)=u_a(q)\cdot\hat r_{\ell a}.
]

## 14.6 Free-volume memory

[
\psi_{\rm fv}(q)=\varphi(q)-\bar\varphi.
]

---

# 15. Orthogonalization of memory coordinates

Memory coordinates must not double-count finite-state indicators.

For each state (A_i),

[
\bar\Psi_i=
\frac{1}{c_i}
\int_{A_i}\Psi(q)\rho_x(\mathrm dq).
]

Define

[
\Psi^\perp(q)=\Psi(q)-\bar\Psi_{s(q)}.
]

The gradient is

[
G_\psi(q)=\nabla_q\Psi^\perp(q)\in\mathbb R^{m\times d}.
]

Within each basin, if (\bar\Psi_{s(q)}) is constant,

[
\nabla_q\Psi^\perp(q)=\nabla_q\Psi(q).
]

---

# 16. Quadrature

For each basin (A_i), choose quadrature points and weights:

[
{(q_{i\ell},w_{i\ell})}_{\ell=1}^{n_i}.
]

Approximate:

[
\int_{A_i} f(q),\mathrm dq
\approx
\sum_{\ell=1}^{n_i}w_{i\ell}f(q_{i\ell}).
]

For transition surfaces or domains, choose points

[
q_{ij,\ell}^{\Sigma}
]

and weights

[
w_{ij,\ell}^{\Sigma}.
]

For transition paths, choose samples

[
\Delta P_{ij,\ell}
]

and weights

[
w_{ij,\ell}^{\rm path}.
]

---

# 17. Mass balance

The model must satisfy conserved component totals.

Define standard concentration

[
C^\circ=1000\ \mathrm{mol,m^{-3}}.
]

Introduce dimensionless chemical potentials

[
\lambda=(\lambda_a)_{a\in\mathcal C}.
]

The basin restricted weight is

[
Z_i(\lambda)=
\int_{A_i}
\exp\left[
-\frac{U_x(q)}{RT}
+
\sum_a\nu_{i,a}\lambda_a
\right],\mathrm dq.
]

Quadrature:

[
Z_i(\lambda)
\approx
\sum_{\ell=1}^{n_i}
w_{i\ell}
\exp\left[
-\frac{U_x(q_{i\ell})}{RT}
+
\sum_a\nu_{i,a}\lambda_a
\right].
]

Basin concentration:

[
c_i(\lambda)=C^\circ Z_i(\lambda).
]

Mass residual:

[
r_a(\lambda)
============

\sum_{i=1}^N\nu_{i,a}c_i(\lambda)-C_a.
]

Solve

[
r_a(\lambda^\star)=0
\qquad
\forall a\in\mathcal C.
]

## 17.1 Jacobian

Since

[
\frac{\partial c_i}{\partial\lambda_b}=\nu_{i,b}c_i,
]

the Jacobian is

[
J_{ab}(\lambda)
===============

# \frac{\partial r_a}{\partial\lambda_b}

\sum_i\nu_{i,a}\nu_{i,b}c_i(\lambda).
]

## 17.2 Newton solve

Set

[
C_{\min}=10^{-30}\ \mathrm{mol,m^{-3}},
]

[
\tau_{\rm mass}=10^{-10}.
]

Initialize:

[
\lambda_a^{(0)}
===============

\log\left(\frac{\max(C_a,C_{\min})}{C^\circ}\right).
]

At iteration (k):

1. Compute (c_i^{(k)}=c_i(\lambda^{(k)})).
2. Compute residual (r^{(k)}).
3. Compute normalized residual

[
\eta^{(k)}
==========

\max_a
\frac{|r_a^{(k)}|}
{\max(C_a,C_{\min})}.
]

4. If

[
\eta^{(k)}<\tau_{\rm mass},
]

stop.

5. Compute (J^{(k)}).
6. Solve

[
J^{(k)}\Delta\lambda^{(k)}=-r^{(k)}.
]

If singular, solve

[
(J^{(k)}+\epsilon_JI)\Delta\lambda^{(k)}=-r^{(k)}
]

with

[
\epsilon_J=10^{-12}\max(1,|J^{(k)}|_2).
]

7. Backtracking:

[
\lambda_{\rm trial}=\lambda^{(k)}+\alpha\Delta\lambda^{(k)}.
]

Start (\alpha=1). Accept if

[
\eta_{\rm trial}<\eta^{(k)}.
]

Otherwise

[
\alpha\leftarrow \alpha/2.
]

Fail if

[
\alpha<2^{-40}.
]

After convergence:

[
\lambda^\star=\lambda.
]

## 17.3 Concentration-density weights

Define

[
W_{i\ell}
=========

C^\circ w_{i\ell}
\exp\left[
-\frac{U_x(q_{i\ell})}{RT}
+
\sum_a\nu_{i,a}\lambda_a^\star
\right].
]

Then

[
c_i=\sum_\ell W_{i\ell}.
]

Define concentration measure

[
\rho_x(\mathrm dq)
==================

C^\circ
\exp\left[
-\frac{U_x(q)}{RT}
+
\sum_a\nu_{s(q),a}\lambda_a^\star
\right]
\mathrm dq.
]

Then

[
c_i=\int_{A_i}\rho_x(\mathrm dq).
]

---

# 18. Committor PDE

For transition edge ((i,j)), define

[
\Omega_{ij}=\Omega_x\setminus(A_i\cup A_j).
]

Define hitting times

[
\tau_i=\inf{t:X_t\in A_i}.
]

The committor is

[
q_{ij}(q)=\Pr_q(\tau_j<\tau_i).
]

It solves

[
L_x q_{ij}=0
\quad\text{in }\Omega_{ij},
]

[
q_{ij}=0
\quad\text{on }\partial A_i,
]

[
q_{ij}=1
\quad\text{on }\partial A_j.
]

## 18.1 Weak form

Find (q_{ij}) with boundary values such that for all test functions (v) vanishing on (\partial A_i\cup\partial A_j),

[
\int_{\Omega_{ij}}
\nabla v(q)^\top D_x(q)\nabla q_{ij}(q),
\rho_x(\mathrm dq)=0.
]

## 18.2 Finite element solve

Choose basis functions

[
{\varphi_m(q)}_{m=1}^M.
]

Represent

[
q_{ij}(q)=q_{ij}^{\rm D}(q)+\sum_{m=1}^M u_m\varphi_m(q),
]

where (q_{ij}^{\rm D}) enforces boundary values.

Stiffness matrix:

[
S_{mn}
======

\int_{\Omega_{ij}}
\nabla\varphi_m(q)^\top D_x(q)\nabla\varphi_n(q),
\rho_x(\mathrm dq).
]

Right-hand side:

[
f_m
===

*

\int_{\Omega_{ij}}
\nabla\varphi_m(q)^\top D_x(q)\nabla q_{ij}^{\rm D}(q),
\rho_x(\mathrm dq).
]

Solve

[
Su=f.
]

Then

[
\nabla q_{ij}(q)
================

\nabla q_{ij}^{\rm D}(q)
+
\sum_m u_m\nabla\varphi_m(q).
]

---

# 19. Symmetric capacity flux (K_{ij})

Compute

[
K_{ij}
======

\int_{\Omega_{ij}}
\nabla q_{ij}(q)^\top D_x(q)\nabla q_{ij}(q),
\rho_x(\mathrm dq).
]

Quadrature:

[
K_{ij}
\approx
\sum_\ell
W_{ij,\ell}^{\Sigma}
g_{ij,\ell}^\top
D_x(q_{ij,\ell}^{\Sigma})
g_{ij,\ell},
]

where

[
g_{ij,\ell}=\nabla q_{ij}(q_{ij,\ell}^{\Sigma}),
]

and

[
W_{ij,\ell}^{\Sigma}
====================

C^\circ w_{ij,\ell}^{\Sigma}
\exp\left[
-\frac{U_x(q_{ij,\ell}^{\Sigma})}{RT}
+
\sum_a\nu_{s(q_{ij,\ell}^{\Sigma}),a}\lambda_a^\star
\right].
]

Set

[
K_{ji}=K_{ij},
\qquad
K_{ii}=0.
]

Units:

[
[K_{ij}]=\mathrm{mol,m^{-3},s^{-1}}.
]

---

# 20. Reversible finite generator (Q)

For (i\ne j),

[
Q_{ij}=\frac{K_{ij}}{c_i}.
]

For diagonal entries,

[
Q_{ii}=-\sum_{j\ne i}Q_{ij}.
]

Required identities:

[
\sum_jQ_{ij}=0,
]

[
c_iQ_{ij}=c_jQ_{ji},
]

[
\sum_i c_iQ_{ij}=0.
]

Units:

[
[Q_{ij}]=\mathrm{s^{-1}}.
]

---

# 21. Transition-path moments (d_{ij},M_{ij})

A transition path from (A_i) to (A_j) carries charge displacement

[
\Delta P_{ij}
=============

P_x(X_{\tau_j})-P_x(X_{\tau_i}).
]

The first moment is

[
d_{ij}
======

\mathbb E[\Delta P_{ij}\mid A_i\to A_j]\in\mathbb R^3.
]

The second moment is

[
M_{ij}
======

\mathbb E[\Delta P_{ij}\Delta P_{ij}^\top\mid A_i\to A_j]
\in\mathbb R^{3\times3}.
]

Reversibility requires

[
d_{ji}=-d_{ij},
]

[
M_{ji}=M_{ij}.
]

## 21.1 Doob-conditioned transition path sampler

The original overdamped drift is

[
b(q)=
\nabla\cdot D_x(q)
------------------

\frac{1}{RT}D_x(q)\nabla U_x(q).
]

The Doob-conditioned drift to hit (A_j) before (A_i) is

[
b^{(ij)}(q)=b(q)+2D_x(q)\nabla\log q_{ij}(q).
]

The conditioned SDE is

[
\mathrm dq_t
============

b^{(ij)}(q_t)\mathrm dt
+
\sqrt{2D_x(q_t)},\mathrm dW_t.
]

Initial distribution on exit surface:

[
p_{ij}^{\rm exit}(q)
====================

\frac{
\rho_x(q),
n_{ij}(q)^\top D_x(q)\nabla q_{ij}(q)
}{
K_{ij}
}.
]

Euler–Maruyama step:

[
q_{n+1}
=======

q_n
+
b^{(ij)}(q_n)\Delta t
+
\sqrt{2D_x(q_n)\Delta t},\eta_n,
\qquad
\eta_n\sim\mathcal N(0,I).
]

Stop when

[
q_t\in A_j.
]

For each sampled path (\ell):

[
\Delta P_{ij,\ell}=P_x(q_{\rm end})-P_x(q_{\rm start}).
]

## 21.2 Weighted estimator

Let weights be (w_{ij,\ell}^{\rm path}). Define

[
S_{ij}^{\rm path}=\sum_\ell w_{ij,\ell}^{\rm path}.
]

Then

[
d_{ij}
======

\frac{1}{S_{ij}^{\rm path}}
\sum_\ell
w_{ij,\ell}^{\rm path}\Delta P_{ij,\ell}.
]

[
M_{ij}
======

\frac{1}{S_{ij}^{\rm path}}
\sum_\ell
w_{ij,\ell}^{\rm path}
\Delta P_{ij,\ell}\Delta P_{ij,\ell}^\top.
]

Set

[
d_{ji}=-d_{ij},
\qquad
M_{ji}=M_{ij},
]

[
d_{ii}=0,
\qquad
M_{ii}=0.
]

## 21.3 Displacement scale validation

Require

[
|\Delta P_{ij,\ell}|\le L_{\max}.
]

For molecular local transitions, a safe default is

[
L_{\max}=10^{-8}\ \mathrm m.
]

A meter-scale displacement is invalid.

---

# 22. Unbounded transport projector (\Pi_i(q))

For each basin (A_i), define

[
\Pi_i(q)\in\mathbb R^{d\times d}.
]

This projects coordinate motion onto unbounded DC charge-transport coordinates.

Examples:

Free ion translation:

[
\Pi_i=I
]

on translational charge coordinates.

Bound internal pair vibration:

[
\Pi_i=0.
]

Neutral pair center-of-mass translation:

[
G_P\Pi_i=0
]

if the pair is exactly neutral and co-moving.

Partner-switch / identity-diffusion coordinate:

[
\Pi_i\neq0
]

only along the unbounded identity coordinate.

The implementation notes emphasize that neutral or bound polarization must not be inserted as an independent charge carrier; multi-center states require charged-center covariance such as (z^\top D_\alpha z), not (z_{\rm net}^2D_{\rm COM}). 

---

# 23. Within-state self-current tensor

For basin (A_i),

[
D_i^{\rm self}
==============

\frac{1}{c_i}
\int_{A_i}
G_P(q)\Pi_i(q)D_x(q)\Pi_i(q)^\top G_P(q)^\top
\rho_x(\mathrm dq).
]

Quadrature:

[
D_i^{\rm self}
\approx
\frac{1}{c_i}
\sum_\ell
W_{i\ell}
G_P(q_{i\ell})
\Pi_i(q_{i\ell})
D_x(q_{i\ell})
\Pi_i(q_{i\ell})^\top
G_P(q_{i\ell})^\top.
]

Require

[
D_i^{\rm self}=D_i^{\rm self\top},
\qquad
D_i^{\rm self}\succeq0.
]

Units:

[
[D_i^{\rm self}]=\mathrm{m^2,s^{-1}}.
]

## 23.1 Multi-center covariance

For a state with center charge vector

[
z_i=(z_1,\ldots,z_m)^\top
]

and center mobility covariance

[
D_i^{\rm centers}\in\mathbb R^{m\times m},
]

the charge mobility is

[
D_Q(i)=z_i^\top D_i^{\rm centers}z_i.
]

For Li–anion:

[
z=(+1,-1)^\top.
]

Then

[
D_Q=D_{\rm Li}+D_{\rm A}-2D_{\rm Li,A}.
]

If

[
D_{\rm Li,A}<0,
]

then

[
-2D_{\rm Li,A}>0
]

and conductivity increases.

This is the mathematical location of bulky-anion / ligand-shell anticorrelation.

---

# 24. Direct diffusivity-density tensor (B)

Define

[
B=
\sum_i c_iD_i^{\rm self}
+
\frac12\sum_i\sum_jK_{ij}M_{ij}.
]

Units:

[
[c_iD_i^{\rm self}]
===================

# \mathrm{mol,m^{-3}}\cdot \mathrm{m^2,s^{-1}}

\mathrm{mol,m^{-1},s^{-1}}.
]

[
[K_{ij}M_{ij}]
==============

# \mathrm{mol,m^{-3},s^{-1}}\cdot \mathrm{m^2}

\mathrm{mol,m^{-1},s^{-1}}.
]

So

[
[B]=\mathrm{mol,m^{-1},s^{-1}}.
]

---

# 25. Finite-state memory / Poisson correction

Define state drift

[
b_i=\sum_jQ_{ij}d_{ij}\in\mathbb R^3.
]

Component:

[
b_{i,a}=\sum_jQ_{ij}d_{ij,a}.
]

For each Cartesian axis (a\in{x,y,z}), solve

[
-Q\chi_a=b_a
]

with gauge

[
\sum_i c_i\chi_{i,a}=0.
]

Because (Q) has a zero eigenvalue, solve the block system

[
\begin{bmatrix}
-Q & \mathbf 1\
c^\top & 0
\end{bmatrix}
\begin{bmatrix}
\chi_a\
\lambda_a^Q
\end{bmatrix}
=============

\begin{bmatrix}
b_a\
0
\end{bmatrix}.
]

The correction tensor is

[
C^Q_{ab}=\sum_i c_i b_{i,a}\chi_{i,b}.
]

Symmetrize numerically:

[
C^Q\leftarrow\frac12(C^Q+C^{Q\top}).
]

---

# 26. Continuous Mori correction

Compute memory gradient

[
G_\psi(q)=\nabla_q\Psi^\perp(q).
]

The Mori Dirichlet matrix is

[
A_{\alpha\gamma}
================

\int
\nabla\psi_\alpha^\perp(q)^\top
D_x(q)
\nabla\psi_\gamma^\perp(q)
\rho_x(\mathrm dq).
]

Quadrature:

[
A
\approx
\sum_i\sum_\ell
W_{i\ell}
G_\psi(q_{i\ell})
D_x(q_{i\ell})
G_\psi(q_{i\ell})^\top.
]

The current coupling is

[
h_{\alpha a}
============

\int
\nabla\psi_\alpha^\perp(q)^\top
D_x(q)
\nabla P_{x,a}(q)
\rho_x(\mathrm dq).
]

Quadrature:

[
h
\approx
\sum_i\sum_\ell
W_{i\ell}
G_\psi(q_{i\ell})
D_x(q_{i\ell})
G_P(q_{i\ell})^\top.
]

Then

[
C^\psi=h^\top A^# h.
]

Symmetrize:

[
C^\psi\leftarrow\frac12(C^\psi+C^{\psi\top}).
]

---

# 27. Projected diffusivity-density tensor

[
D_{\rm proj}=B-C^Q-C^\psi.
]

Symmetrize:

[
D_{\rm proj}\leftarrow\frac12(D_{\rm proj}+D_{\rm proj}^\top).
]

If numerical roundoff causes tiny negative eigenvalues, eigendecompose

[
D_{\rm proj}=V\operatorname{diag}(\eta_1,\eta_2,\eta_3)V^\top.
]

If

[
\min_r\eta_r\ge -\tau_D\max_r|\eta_r|
]

with e.g.

[
\tau_D=10^{-10},
]

then set

[
\eta_r^+=\max(\eta_r,0)
]

and reconstruct

[
D_{\rm proj}^+=V\operatorname{diag}(\eta_1^+,\eta_2^+,\eta_3^+)V^\top.
]

If the negative eigenvalue is larger than tolerance, fail.

---

# 28. Final scalar conductivity

The isotropic diffusivity-density is

[
D_{\rm iso}=\frac13\operatorname{Tr}(D_{\rm proj}).
]

Using formal-charge displacement units, conductivity is

[
\sigma_{\rm S/m}
================

\frac{F^2}{RT}D_{\rm iso}.
]

Therefore

[
\boxed{
\sigma_{\rm S/m}
================

\frac{F^2}{3RT}
\operatorname{Tr}
\left[
\sum_i c_iD_i^{\rm self}
+
\frac12\sum_{i,j}K_{ij}M_{ij}
-----------------------------

## C^Q

h^\top A^#h
\right].
}
]

And

[
\boxed{
\sigma_{\rm mS/cm}
==================

10\frac{F^2}{3RT}
\operatorname{Tr}
\left[
\sum_i c_iD_i^{\rm self}
+
\frac12\sum_{i,j}K_{ij}M_{ij}
-----------------------------

## C^Q

h^\top A^#h
\right].
}
]

This is the final scalar output.

---

# 29. Full algorithm

Input:

[
x=(T,{C_a},{\phi_s},\mathcal L).
]

Output:

[
\sigma_{\rm mS/cm}(x).
]

Algorithm:

1. Build molecule/site counts from (C_a,\phi_s,\mathcal L).
2. Build coordinate space (\Omega_x) and coordinate vector (q).
3. Build (P_x(q)) and (G_P(q)) from charge sites.
4. Build (U_x(q)) from topology and force-field terms.
5. Build (D_x(q)) from resistance/mobility library.
6. Build basins (A_i) from physical coordinates.
7. Build stoichiometry vectors (\nu_i).
8. Generate basin quadrature ((q_{i\ell},w_{i\ell})).
9. Solve chemical potentials (\lambda^\star) by damped Newton.
10. Compute weights (W_{i\ell}).
11. Compute populations (c_i).
12. For each transition edge ((i,j)), solve committor PDE.
13. Compute capacity flux (K_{ij}).
14. Build reversible generator (Q).
15. Sample Doob-conditioned local transition paths.
16. Compute (d_{ij},M_{ij}).
17. Define unbounded projectors (\Pi_i).
18. Compute (D_i^{\rm self}).
19. Compute direct tensor (B).
20. Compute state drift (b_i).
21. Solve finite-state Poisson equations.
22. Compute (C^Q).
23. Build memory coordinates (\Psi).
24. Orthogonalize to (\Psi^\perp).
25. Compute (A,h).
26. Compute (C^\psi=h^\top A^#h).
27. Compute (D_{\rm proj}=B-C^Q-C^\psi).
28. Return

[
\sigma_{\rm mS/cm}=10\frac{F^2}{3RT}\operatorname{Tr}(D_{\rm proj}).
]

---

# 30. Basis refinement without production GK/EH

Let current memory basis be

[
\Phi=(\phi_1,\ldots,\phi_m).
]

For candidate missing coordinate (g(q)), compute

[
A_{gg}=
\int
\nabla g^\top D_x\nabla g,
\rho_x(\mathrm dq),
]

[
A_{g\Phi}=
\int
\nabla g^\top D_x\nabla\Phi,
\rho_x(\mathrm dq),
]

[
h_g=
\int
\nabla g^\top D_xG_P^\top,
\rho_x(\mathrm dq).
]

Let

[
A_{\Phi\Phi}=
\int
\nabla\Phi^\top D_x\nabla\Phi,
\rho_x(\mathrm dq),
]

[
h_\Phi=
\int
\nabla\Phi^\top D_xG_P^\top,
\rho_x(\mathrm dq).
]

Remove current-basis projection:

[
\tilde h_g=
h_g-A_{g\Phi}A_{\Phi\Phi}^#h_\Phi,
]

[
\tilde A_{gg}=
A_{gg}-A_{g\Phi}A_{\Phi\Phi}^#A_{\Phi g}.
]

Score:

[
S(g)=\frac{|\tilde h_g|*2^2}{\tilde A*{gg}}.
]

Add

[
g^\star=\arg\max_gS(g)
]

if

[
S(g^\star)>\tau_{\rm basis}.
]

Stop when

[
\max_gS(g)<\tau_{\rm basis}
]

and

[
|\sigma_{n+1}-\sigma_n|<\tau_\sigma.
]

This basis-refinement scheme is in the retrieved model notes and is explicitly defined without a production Green–Kubo or Einstein–Helfand calculation. 

---

# 31. Convergence proof

Let exact conductivity be

[
\sigma_{\rm GK}
===============

\frac{F^2}{3RT}
\sum_{a=1}^3
\langle J_a,(-L_x)^#J_a\rangle_{\mu_x}.
]

Let (V_n) be the span of all state indicators and memory coordinates at level (n).

Let (u_a^\star) solve

[
-L_xu_a^\star=J_a.
]

Let (u_{n,a}\in V_n) solve the Galerkin projection:

[
\langle v,(-L_x)u_{n,a}\rangle_{\mu_x}
======================================

\langle v,J_a\rangle_{\mu_x}
\qquad
\forall v\in V_n.
]

Define energy norm

[
|f|*E^2=
\langle f,(-L_x)f\rangle*{\mu_x}.
]

Then

[
\sigma_{\rm GK}-\sigma_n
========================

\frac{F^2}{3RT}
\sum_{a=1}^3
|u_a^\star-u_{n,a}|_E^2.
]

If

[
\overline{\bigcup_n V_n}^{|\cdot|_E}
]

contains every (u_a^\star), then

[
\sigma_n\to\sigma_{\rm GK}.
]

The retrieved proof note states this same energy-norm convergence form for the projected model. 

---

# 32. Conductivity-effect identification

An effect is any composition, coordinate, or parameter perturbation

[
\theta
]

that changes at least one primitive:

[
c,\ K,\ Q,\ d,\ M,\ D^{\rm self},\ A,\ h.
]

For small perturbation (\delta\theta), compute central difference:

[
\partial_\theta X
\approx
\frac{
X(\theta+\delta\theta)-X(\theta-\delta\theta)
}{
2\delta\theta
}
]

for

[
X\in{c,K,d,M,D^{\rm self},A,h,\sigma}.
]

Primitive ownership scores:

[
S_c(\theta)=|\partial_\theta c|_2,
]

[
S_K(\theta)=|\partial_\theta K|_F,
]

[
S_{dM}(\theta)=
|\partial_\theta d|*F+|\partial*\theta M|_F,
]

[
S_D(\theta)=|\partial_\theta D^{\rm self}|_F,
]

[
S_{Ah}(\theta)=|\partial_\theta A|*F+|\partial*\theta h|_F.
]

Assign the effect to the largest nonzero score.

## 32.1 Effect ledger

| Effect                                   | Primitive location                 | Computation                                                         |                                  |
| ---------------------------------------- | ---------------------------------- | ------------------------------------------------------------------- | -------------------------------- |
| Free-ion fraction                        | (c_i)                              | (\partial c_i/\partial\theta)                                       |                                  |
| Ion association                          | (c_i,K_{ij})                       | (\partial c,\partial K)                                             |                                  |
| SSIP/CIP balance                         | (c_i,K_{ij},D_i^{\rm self})        | (\partial c,\partial K,\partial D)                                  |                                  |
| Aggregate formation                      | (c_i,K_{ij},M_{ij},D_i^{\rm self}) | (\partial c,\partial K,\partial M,\partial D)                       |                                  |
| Neutral ligand coordination              | (c_i,K_{ij},A,h)                   | (\partial c,\partial K,\partial A,\partial h)                       |                                  |
| Solvation competition                    | (U_x\to c_i,K_{ij})                | (\partial U\to\partial c,\partial K)                                |                                  |
| Coulomb association                      | (U_x\to c_i,K_{ij})                | (\partial U\to\partial c,\partial K)                                |                                  |
| Born/desolvation                         | (U_x\to c_i,K_{ij})                | (\partial U\to\partial c,\partial K)                                |                                  |
| Activity/nonideality                     | (U_x\to c_i,K_{ij})                | (\partial U\to\partial c,\partial K)                                |                                  |
| Viscosity                                | (D_x\to D_i^{\rm self},K_{ij},A)   | (\partial D\to\partial D^{\rm self},\partial K,\partial A)          |                                  |
| Hydrodynamic size                        | (D_x\to D_i^{\rm self},K_{ij})     | (\partial D)                                                        |                                  |
| Anion shape/denticity                    | (U_x,D_x,K_{ij},A,h)               | (\partial U,\partial D,\partial K,\partial A,\partial h)            |                                  |
| Charge-cloud radius                      | (U_x,D_x,A,h)                      | (\partial U,\partial D,\partial A,\partial h)                       |                                  |
| Local crowding/free volume               | (U_x,D_x,K_{ij},A,h)               | (\partial U,\partial D,\partial K,\partial A,\partial h)            |                                  |
| Li–anion co-motion                       | (D_i^{\rm self},M_{ij})            | (D_{\rm Li,A}>0)                                                    |                                  |
| Li–anion anti-correlation                | (D_i^{\rm self},M_{ij})            | (D_{\rm Li,A}<0)                                                    |                                  |
| Partner switching                        | (K_{ij},d_{ij},M_{ij},A,h)         | (\partial K,\partial d,\partial M,\partial A,\partial h)            |                                  |
| Identity diffusion                       | (K_{ij},d_{ij},M_{ij})             | (\partial K,\partial d,\partial M)                                  |                                  |
| Structural hopping                       | (K_{ij},d_{ij},M_{ij})             | (\partial K,\partial d,\partial M)                                  |                                  |
| Bound neutral polarization               | (A,h)                              | (\partial A,\partial h)                                             |                                  |
| Cage/backjump                            | (C^Q) or (A,h)                     | (\partial C^Q,\partial A,\partial h)                                |                                  |
| Ion-atmosphere relaxation                | (A,h)                              | (\partial A,\partial h)                                             |                                  |
| Electrophoretic drag                     | (D_x) or (A,h)                     | (\partial D,\partial A,\partial h)                                  |                                  |
| Debye screening                          | (U_x,D_x,A,h)                      | (\partial U,\partial D,\partial A,\partial h)                       |                                  |
| Solvent-shell residence                  | (c_i,K_{ij},A,h)                   | (\partial c,\partial K,\partial A,\partial h)                       |                                  |
| Ligand-shell residence                   | (c_i,K_{ij},A,h)                   | (\partial c,\partial K,\partial A,\partial h)                       |                                  |
| Temperature                              | (RT,U_x,D_x,K_{ij},F^2/(RT))       | (\partial_T\sigma)                                                  |                                  |
| Dielectric response                      | (U_x,D_x,A,h)                      | (\partial U,\partial D,\partial A,\partial h)                       |                                  |
| Salt concentration                       | (c_i,U_x,D_x,K_{ij},A,h)           | (\partial c,\partial U,\partial D,\partial K,\partial A,\partial h) |                                  |
| Additive microviscosity                  | (D_x,K_{ij})                       | (\partial D,\partial K)                                             |                                  |
| Bulky-anion/ligand conductivity increase | (D_i^{\rm self}) or (M_{ij})       | (D_{\rm Li,A                                                        | i}<0) or partner-switch (M_{ij}) |

The retrieved effect-checklist says the same: an effect is real only if it is identifiable in (c_i), (K_{ij}/Q_{ij}), (d_{ij}/M_{ij}), (D_i/R_i), or (A,h); a final scalar correction is not a physical effect. 

---

# 33. Validation rules

A numerical implementation must fail if any of the following occur:

[
c_i\le0,
]

[
K_{ij}<0,
]

[
K_{ij}\ne K_{ji},
]

[
\sum_jQ_{ij}\ne0,
]

[
c_iQ_{ij}\ne c_jQ_{ji},
]

[
d_{ji}\ne -d_{ij},
]

[
M_{ji}\ne M_{ij},
]

[
D_i^{\rm self}\not\succeq0,
]

[
A\not\succeq0,
]

[
|\Delta P_{ij,\ell}|>L_{\max}.
]

The implementation must also fail if a recipe-only row is converted into a fake generator without the physical library (\mathcal L).

---

# 34. Complete final formula

[
\boxed{
\sigma_{\rm mS/cm}(x)
=====================

10\frac{F^2}{3RT}
\operatorname{Tr}
\left[
\sum_i c_iD_i^{\rm self}
+
\frac12\sum_{i,j}K_{ij}M_{ij}
-----------------------------

## C^Q

h^\top A^#h
\right].
}
]

with

[
C^Q_{ab}=\sum_i c_i b_{i,a}\chi_{i,b},
]

[
b_i=\sum_jQ_{ij}d_{ij},
]

[
-Q\chi_a=b_a,
\qquad
\sum_i c_i\chi_{i,a}=0,
]

[
A_{\alpha\gamma}
================

\int
\nabla\psi_\alpha^\perp{}^\top
D_x
\nabla\psi_\gamma^\perp
,\rho_x(\mathrm dq),
]

[
h_{\alpha a}
============

\int
\nabla\psi_\alpha^\perp{}^\top
D_x
\nabla P_a
,\rho_x(\mathrm dq).
]

That is the full model.
